\originalchapter{生物状态与随机演化}
\originalsection{分子状态和时间尺度}
活细胞中的反应物不是连续液体中的一个抽象浓度，而是一批有限、可以逐个计数的分子。若某种调控蛋白只有几十个拷贝，一次结合就会明显改变可用分子的数目。相同初始条件的两次观测也不必得到相同轨迹。统计模型要预测的是轨迹的分布，而非宣称每条轨迹都满足平均方程。

在室温附近，热能 $\kb T$ 约为 $4.1\,\mathrm{pN\,nm}$。将分子大小取为 $a$，扩散系数取为 $D$，则扩散跨越此长度所需时间约为 $a^2/D$；化学反应还有独立的等待时间。只有在位置混合比反应快、外部条件在观测期间近似不变时，才可把空间自由度消去，用一个常数反应率描述状态转移。

设 $p_i(t)$ 是状态 $i$ 的概率，$W_{ij}\geq0$ 是从 $j$ 到 $i$ 的每单位时间跃迁率。小时间 $\Delta t$ 内从 $j$ 到 $i$ 的转移概率为 $W_{ij}\Delta t+o(\Delta t)$；两次以上转移属于高阶项。收集流入与流出便得到主方程
\begin{equation}\label{eq:master}
\dot p_i=\sum_{j\ne i}(W_{ij}p_j-W_{ji}p_i).
\end{equation}
对 $i$ 求和时，每条边的流入与流出抵消，所以归一化保持。对角元必须是该状态全部流出率的负值，不能直接填一个独立参数。

\begin{example}
单个分子在活性态 $A$ 与非活性态 $B$ 间以 $A\xrightarrow{a}B$、$B\xrightarrow{b}A$ 转换。初始处于 $A$，求活性概率和弛豫时间。
\end{example}
\begin{solution}
$\dot p_A=-ap_A+b(1-p_A)$，故 $p_A(t)=b/(a+b)+a\e^{-(a+b)t}/(a+b)$。平衡由两种方向的概率流相等确定，弛豫时间为 $(a+b)^{-1}$。这不是平均停留时间：在 $A$ 中单次停留时间的均值是 $a^{-1}$。
\end{solution}

若 $N$ 个分子独立转换，活性分子数的概率是二项分布。独立性同时决定方差 $Np_A(1-p_A)$；有共享酶、竞争结合或协同反应时，不可继续套用这一方差。对 $A+B\to C$ 之类双分子反应，状态转移率还随两种反应物的拷贝数变化。

\originalsection{连续极限与遗传漂变}
随机繁殖提供另一个离散模型。取一个固定大小为 $N$ 的单倍体种群，当前等位基因频率为 $x$。下一代独立抽样，则拷贝数 $m\sim\mathrm{Bin}(N,x)$，$\Delta x=m/N-x$ 满足 $\avg{\Delta x}=0$ 和 $\avg{(\Delta x)^2}=x(1-x)/N$。没有选择优势仍会出现频率漂移，因为一次有限抽样不等于其期望。

将转移核对小频率增量展开，并以一代为时间单位，二阶连续近似给出
\begin{equation}\label{eq:fp}
\partial_t p=-\partial_x[A(x)p]+\tfrac12\partial_x^2[B(x)p],\qquad B(x)=\frac{x(1-x)}N.
\end{equation}
这里 $A$ 是每代平均频率变化，$B$ 是每代频率变化的方差，而不是通常写成 $2D$ 之后的 $D$。中性模型 $A=0$。$x=0,1$ 是吸收态：失去的等位基因不能由纯抽样重新产生。连续密度的积分会下降，失去的质量转移到边界的概率原子，不能把它误判为概率不守恒。

设 $u(x)$ 为最终固定到 $1$ 的概率。从第一步条件期望 $u(x)=\avg{u(x+\Delta x)}$ 展开，得到后向方程 $Au'+Bu''/2=0$。中性情形与边界 $u(0)=0,u(1)=1$ 给出 $u(x)=x$。平均吸收时间 $T$ 的第一步式为 $T(x)=1+\avg{T(x+\Delta x)}$，因此
\[
\frac{x(1-x)}{2N}T''=-1,\qquad
T(x)=-2N[x\ln x+(1-x)\ln(1-x)].
\]
积分常数由两个吸收边界确定，$x\ln x\to0$ 保证极限有限。若是二倍体且 $N$ 指个体数，应把基因拷贝数改为 $2N$；不同教材中的因子二常来自这一约定。

\originalsection{突变、选择与稳态}
对称突变率 $\mu$ 加入漂移 $A=\mu(1-2x)$。零概率流稳态满足 $Ap-(Bp)'/2=0$，积分后为 $p_*(x)\propto[x(1-x)]^{2N\mu-1}$。当 $2N\mu<1$ 时密度在两端高，但只要 $\mu>0$，边界奇性仍可积；它反映种群常常几乎单型，而不意味着已经永久吸收。选择优势 $s$ 又加入 $s x(1-x)$，稳态多出因子 $\e^{2Nsx}$。

用数据反推 $N\mu$ 时必须考虑取样数量、低频变异漏检以及人群历史。将两个群体的直方图差异直接解释为一种生物机制，超出了这个平稳、恒定种群模型。8.592J 的相关公开数据可用于讨论这些限制，本卷没有把未经拟合的数据比较写成已经取得的实验结论。
